\needspace{6\baselineskip}
\section{Reproducibility}\label{app:reproducibility}
Archive and Lean sources:
\href{https://doi.org/10.5281/zenodo.24254845}{doi:10.5281/zenodo.24254845}.

To produce \texttt{main.pdf}, run \texttt{pdflatex main.tex} three
times from the arXiv source root containing \texttt{main.tex},
\texttt{references.tex}, \texttt{sections/} and \texttt{figures/}.
The tested engine was MiKTeX-pdfTeX 4.23 (MiKTeX 25.12).

Fix the universal machine and compiler of Sections~\ref{intro:question} and~\ref{s:model}
and the classical imports in Appendix~\ref{app:imports}. For Theorem~\ref{thm:main},
start with the controller of Section~\ref{sec:controller}, use the
compiler and constant choices of Section~\ref{sec:compilation}, and
construct the quantum schedule of Section~\ref{sec:quantum}. For
Theorem~\ref{thm:exponential}, use the affine layout and static chart of
Section~\ref{sec:exponential}, then the parameter order of
Section~\ref{v:parameters}. For Theorem~\ref{thm:four-dimensional}, use the folded geometry of
Section~\ref{f4:geometry-section} and the parameter order of
Section~\ref{f4:constants}. The recipe determines the circuits once
the choices allowed by the stated construction bounds are fixed.
Every realization satisfying those bounds obeys the corresponding theorem, with its own constants and windows.

From a circuit choice and label \(n\), compute the physical rate window
by the rational inner-endpoint prescriptions of
Sections~\ref{s:parameters} and~\ref{v:parameters} or Lemma~\ref{f4:encoder},
as appropriate. From \((n,L)\), the finite initialization searches of
Lemmas~\ref{c:initialization}, \ref{v:encoder} and~\ref{f4:encoder}
give the encoder, using the exact code procedures in Lemma~\ref{d:code}
for the concatenated code. For the coexistence circuit, start with
the wrapped machine, folded track and closed-code track in the proof
of Proposition~\ref{p:coexistence}. That proof specifies the rational
inner physical windows, the closed-code initialization on \(I_0\)
and the folded encoder used for the tail and for linear certification.

For a readout at \((n,L,t)\), Proposition~\ref{d:storage} and the
proof of Theorem~\ref{f4:detailed}(i) prescribe noiseless completion
with the actual in-flight work followed by the fixed bottom-up recovery.
Proposition~\ref{v:certification}(ii) fixes the intended-complex
decoder, including lexicographic box, lift and filling choices and
a fixed failure correction. To query the prescribed fidelity at a
rational rate \(q\), use the finite polynomial calculation of
Proposition~\ref{a:rational} and the six-dimensional adaptation in
Proposition~\ref{v:certification}(ii).

These reconstructions are algorithms specified in the paper. The archive
provides no programs other than the Lean library of
Appendix~\ref{app:lean}, which checks the two results listed there.
